% Standalone problem document, generated from the adjacent README.md.
% Compile with XeLaTeX. Mathematical content is maintained in Markdown.
\documentclass[11pt,a4paper]{article}
\usepackage[fontsize=10.5pt]{scrextend}
\usepackage[margin=22mm,headheight=15pt,headsep=9mm,footskip=11mm]{geometry}
\usepackage{fontspec}
\setmainfont{texgyrepagella}[Extension=.otf,UprightFont=*-regular,BoldFont=*-bold,ItalicFont=*-italic,BoldItalicFont=*-bolditalic]
\setsansfont{texgyreheros}[Extension=.otf,UprightFont=*-regular,BoldFont=*-bold,ItalicFont=*-italic,BoldItalicFont=*-bolditalic]
\setmonofont{texgyrecursor}[Extension=.otf,UprightFont=*-regular,BoldFont=*-bold,ItalicFont=*-italic,BoldItalicFont=*-bolditalic,Scale=MatchLowercase]
\usepackage{amsmath,amssymb}
\usepackage{unicode-math}
\setmathfont{texgyrepagella-math.otf}
\usepackage{xcolor,microtype,enumitem,needspace,fancyhdr}
\usepackage{bookmark}
\definecolor{ink}{HTML}{17344B}
\definecolor{accent}{HTML}{197D80}
\definecolor{muted}{HTML}{596773}
\hypersetup{colorlinks=true,linkcolor=accent,urlcolor=accent,pdftitle={MI-30 - Product
inequality for disjoint principal minors of Wishart
matrices},pdfauthor={Open Problems in Numerical Linear Algebra}}
\urlstyle{same}
\setlength{\parindent}{0pt}
\setlength{\parskip}{5pt plus 1pt minus 1pt}
\linespread{1.04}
\setlength{\emergencystretch}{3em}
\widowpenalty=10000
\clubpenalty=10000
\displaywidowpenalty=10000
\allowdisplaybreaks[1]
\setlist{nosep,leftmargin=1.5em}
\providecommand{\tightlist}{\setlength{\itemsep}{0pt}\setlength{\parskip}{0pt}}
\providecommand{\pandocbounded}[1]{#1}
\pagestyle{fancy}
\fancyhf{}
\fancyhead[L]{\sffamily\footnotesize\color{muted}OPEN PROBLEMS IN NUMERICAL LINEAR ALGEBRA}
\fancyhead[R]{\sffamily\bfseries\footnotesize\color{accent}MI-30}
\fancyfoot[L]{\parbox[b]{0.9\textwidth}{\sffamily\fontsize{8}{10}\selectfont\color{muted}Editorial ratings; a bounded search cannot certify that a problem remains unsolved.\\Literature check: 2026-09-11}}
\fancyfoot[R]{\sffamily\footnotesize\color{muted}\thepage}
\renewcommand{\headrulewidth}{0.3pt}
\renewcommand{\headrule}{\hbox to\headwidth{\color{accent}\leaders\hrule height \headrulewidth\hfill}}
\makeatletter
\renewcommand{\section}{\Needspace{5\baselineskip}\@startsection{section}{1}{0pt}{12pt plus 2pt minus 2pt}{4pt}{\sffamily\bfseries\large\color{ink}}}
\renewcommand{\subsection}{\Needspace{4\baselineskip}\@startsection{subsection}{2}{0pt}{9pt plus 1pt minus 1pt}{3pt}{\sffamily\bfseries\normalsize\color{ink}}}
\makeatother
\setcounter{secnumdepth}{0}
\begin{document}
{\sffamily\small\color{muted}Matrix inequalities and norms\par}
\vspace{3pt}
{\raggedright\hyphenpenalty=10000\exhyphenpenalty=10000\sffamily\bfseries\fontsize{21}{24}\selectfont\color{ink}Product
inequality for disjoint principal minors of Wishart matrices\par}
\vspace{7pt}
{\sffamily\small\textbf{Difficulty:} challenging\quad\textbf{Importance:} interesting
to specialist\par}
{\sffamily\small\bfseries\color{accent}Status: Partially resolved\par}
\vspace{4pt}
{\color{accent}\hrule height 0.8pt}
\vspace{6pt}
\textbf{Topic:} Random Gram determinants and Gaussian product
inequalities

\textbf{Rating rationale:} The two-block result does not yield the
general joint inequality, which couples arbitrarily many determinant
powers. The main audience studies random Gram matrices, covariance
determinants, and Gaussian volume inequalities.

\subsection{Statement}\label{statement}

Let \(d,p_1,\ldots,p_d\) be positive integers, set
\(p=\sum_{i=1}^d p_i\), and choose a real number \(\alpha>p-1\) and a
real symmetric positive definite matrix
\(\Sigma\in\mathbb R^{p\times p}\).

Let \(X\sim W_p(\alpha,\Sigma)\) mean that \(X\) has density
proportional to

\[
(\det X)^{(\alpha-p-1)/2}
\exp\!\left[-\tfrac12\operatorname{tr}(\Sigma^{-1}X)\right]
\]

on the cone of real symmetric positive definite matrices, with respect
to Lebesgue measure on their independent entries. Partition
\(X=(X_{ij})_{i,j=1}^d\) so that \(X_{ij}\) has size \(p_i\times p_j\).

\textbf{Conjecture (Genest--Ouimet--Richards).} For every such choice
and every \(\nu_1,\ldots,\nu_d\ge0\),

\[
\mathbb E\left[\prod_{i=1}^d(\det X_{ii})^{\nu_i}\right]
\ge
\prod_{i=1}^d\mathbb E\left[(\det X_{ii})^{\nu_i}\right].
\]

All displayed expectations are finite. Degrees of freedom and exponents
need not be integers. All block counts and block sizes belong to this
one conjecture.

\subsection{Known cases and numerical
significance}\label{known-cases-and-numerical-significance}

The conjecture is proved for \(d=2\). Block-diagonal \(\Sigma\) makes
the diagonal blocks independent and gives equality. Results for negative
determinant powers address a different exponent regime.

For integer \(\alpha\ge p\), one may write \(X=Y^{\mathsf T}Y\), where
\(Y\) has \(\alpha\) independent Gaussian rows with covariance
\(\Sigma\). Then \(\det X_{ii}\) is the squared volume generated by the
columns in block \(i\). The problem compares joint volume moments with
the corresponding product of marginal moments, linking it to random Gram
matrices and volume methods.

\newpage
\subsection{References and status
check}\label{references-and-status-check}

\begin{itemize}
\tightlist
\item
  C. Genest, F. Ouimet and D. Richards, \emph{On the Gaussian product
  inequality conjecture for disjoint principal minors of Wishart random
  matrices}, Electronic Journal of Probability 29 (2024), paper 166.
  \href{https://doi.org/10.1214/24-EJP1222}{DOI};
  \href{https://arxiv.org/html/2311.00202v3}{current arXiv v3}, dated
  2024-10-01. Conjecture 1.1, equation (6), states the target;
  Definition 2.4 defines the distribution; Remarks 1.2--1.3 record the
  two-block case and moment finiteness.
\item
  C. Genest, F. Ouimet and D. Richards, \emph{An explicit Wishart moment
  formula for the product of two disjoint principal minors}, Proceedings
  of the AMS 153 (2025), 1299--1311.
  \href{https://doi.org/10.1090/proc/17077}{DOI};
  \href{https://arxiv.org/html/2409.14512}{full text}. This proves the
  two-block case.
\item
  R. E. Gaunt and F. Ouimet, \emph{On Wilks' problem: Exact recursive
  formulas via Stein's method for the joint moments of disjoint
  principal minors of Wishart random matrices},
  \href{https://arxiv.org/html/2607.03984}{arXiv:2607.03984} (2026). Its
  moment recursions do not assert the displayed inequality for all real
  exponents.
\item
  F. Ouimet and D. Greaves, \emph{A proof of the strong Gaussian product
  inequality conjecture} (2026),
  \href{https://www.researchgate.net/profile/Frederic-Ouimet/publication/410720385_A_proof_of_the_strong_Gaussian_product_inequality_conjecture/links/6a67b848f475b23f400ec34e/A-proof-of-the-strong-Gaussian-product-inequality-conjecture.pdf}{author-posted
  manuscript}, Theorem 2.1. This claims a scalar Gaussian-coordinate
  product inequality; the introduction separately identifies Wishart
  extensions. Its stated result does not cover arbitrary determinant
  blocks. This entry does not assess that manuscript's proof.
\end{itemize}

On 2026-09-11, checked the source's current version, these later papers
and \href{https://sites.google.com/site/fouimet26/research}{Ouimet's
publication list}, and searched the exact title, Wishart product
inequalities, Conjecture 1.1, and 2025--2026 proof/counterexample
combinations. No full resolution of the displayed determinant-block
target was found. This is a bounded status check.
\end{document}
